\section{Validation and performance}
\label{sec:validation}

Two questions are often run together and answered with the same number. Does
\phynetpy\ compute what it claims to compute? And how fast is it? The first is
a correctness question with a right answer; the second is a measurement that
depends on hardware, data, and what each implementation was asked to do. This
section keeps them apart.

All measurements below were taken on the platform recorded in
Section~\ref{sec:reproducibility} and are regenerated by scripts in the
repository; no figure in this section was entered by hand.

\subsection{Agreement with PhyloNet}

The strongest correctness evidence available for a reimplementation is
numerical agreement with the implementation it reimplements, on inputs chosen
to be difficult. \phynetpy\ ships a harness that compiles a Java driver
invoking PhyloNet's own classes -- \code{GeneTreeBrSpeciesNetDistribution} for
the multispecies-network-coalescent density, and its BEAGLE-backed
\cite{Ayres2012BEAGLE} \code{UltrametricTree} for the Felsenstein likelihood --
and feeds both implementations from one shared specification file, so a
disagreement is a real discrepancy rather than a difference in test setup.

Fourteen cases are specified, chosen to probe the places where two
implementations of the same density diverge: multiple alleles per species, GTR
with skewed base frequencies and asymmetric exchange rates, stacked
reticulations, boundary population sizes, embeddings of probability zero (where
both sides must agree on $-\infty$), IUPAC ambiguity codes and gaps,
branch-length saturation, fully constant alignments, and a five-taxon
caterpillar. The two quantities compared are
$\log P(g \mid \Psi)$ and $\log P(S \mid g)$.

All cases agree to within $10^{-5}$, and the $-\infty$ case matches exactly. The
suite runs as 23 pytest cases and skips cleanly when no PhyloNet jar or BEAGLE
installation is present, so it does not burden contributors who have neither.

This is a statement about two likelihood functions, not about two searches.
Agreement on a fixed network says nothing about whether the two searches visit
the same topologies, and Section~\ref{sec:headtohead} shows that on reticulate
data they do not always find the same answer.

\subsection{Other correctness evidence}

\begin{table}[t]
  \centering
  \small
  % Generated by benchmarks/make_artifacts.py -- do not edit by hand.
\begin{tabular}{lr}
\toprule
Area & Tests \\
\midrule
Substitution models (\texttt{GTR.py}) & 432 \\
Network comparison and decomposition & 153 \\
Network representation and topology moves & 163 \\
Inference API and dispatch registry & 88 \\
PhyloNet numerical cross-check & 23 \\
Memory and mutation-cost bounds & 9 \\
\midrule
\textbf{Total collected} & \textbf{1130} \\
Passing in the default run & 1125 \\
Deselected as \texttt{slow} & 5 \\
Failing & 0 \\
\midrule
\multicolumn{2}{l}{Wall time, default run: 46\,s} \\
\multicolumn{2}{l}{CI Python versions: 3.9, 3.10, 3.11, 3.12, 3.13, 3.14} \\
\multicolumn{2}{l}{CI jobs: tests, windows, package, slow} \\
\bottomrule
\end{tabular}

  \caption{Test counts by area, and the continuous-integration configuration.
  Counts are collected by \code{benchmarks/bench_testsuite.py}, which runs the
  suite and parses the result rather than restating it. The skipped count is
  environment dependent: the PhyloNet cross-check runs only where a jar and
  BEAGLE are installed, and it did run for the figures reported here.}
  \label{tab:testing}
\end{table}

Table~\ref{tab:testing} gives the test suite by area. Three of its entries are
doing more than counting.

The 432 substitution-model tests exist because an audit of
\code{GTR.py} found that its rate matrix was not time reversible for
non-uniform base frequencies, that $e^{Qt}$ was being computed with a symmetric
eigensolver on a non-symmetric matrix, and that three of the eight models could
not be constructed at all. Every model is now checked against
\code{scipy.linalg.expm} \cite{Virtanen2020SciPy}; the largest absolute
disagreement across the eight is $4.4 \times 10^{-16}$
(Table~\ref{tab:substitution}). Reporting this is uncomfortable but
informative, and the reason no published result changes is worth stating
plainly: nothing on the inference path used those classes. The sequence
likelihood has its own, separately correct implementation, so the audit
cross-validates two implementations that now agree exactly rather than
correcting any analysis.

The incremental pseudo-likelihood scorer -- which reuses a cached dynamic
program when only branch lengths or inheritance probabilities change -- is
checked against a full rebuild rather than against expectation, which is what
makes the optimisation safe to leave enabled.

The nine tests on memory and mutation cost are regression bounds rather than
measurements: they pair timing and allocation limits with structural invariant
checks on incidence bookkeeping, so a change that degrades edge-mutation cost or
corrupts the indices fails the suite instead of being noticed later.

\subsection{Head-to-head against PhyloNet}
\label{sec:headtohead}

\begin{table}[t]
  \centering
  \scriptsize
  % Generated by benchmarks/make_artifacts.py -- do not edit by hand.
\begin{tabular}{rrrrrrrrrr}
\toprule
& & & & \multicolumn{3}{c}{Mean runtime (s)} & \multicolumn{3}{c}{Replicates recovering truth} \\
\cmidrule(lr){5-7}\cmidrule(lr){8-10}
Taxa & Genes & $r$ & Conc. & \textsc{PNPy} & \textsc{PNPy} & PhyloNet & \textsc{PNPy} & \textsc{PNPy} & PhyloNet \\
& & & & (default) & (\texttt{phylonet}) & & (default) & (\texttt{phylonet}) & \\
\midrule
4 & 25 & 0 & 40\% & 0.11 & 0.08 & 0.11 & 2/3 & 2/3 & 2/3 \\
4 & 100 & 0 & 46\% & 0.09 & 0.06 & 0.15 & 1/3 & 1/3 & 2/3 \\
5 & 50 & 0 & 34\% & 0.12 & 0.12 & 0.45 & 1/3 & 1/3 & 2/3 \\
5 & 200 & 0 & 38\% & 0.12 & 0.12 & 0.41 & 1/3 & 1/3 & 2/3 \\
6 & 100 & 0 & 22\% & 0.19 & 0.15 & 0.70 & 0/3 & 0/3 & 2/3 \\
6 & 400 & 0 & 26\% & 0.17 & 0.12 & 0.73 & 0/3 & 0/3 & 1/3 \\
8 & 200 & 0 & 15\% & 0.41 & 0.46 & 2.38 & 1/3 & 1/3 & 3/3 \\
5 & 100 & 1 & 0\% & 1.11 & 5.48 & 0.86 & 2/3 & 1/3 & 1/3 \\
6 & 200 & 1 & 16\% & 1.08 & 0.93 & 1.48 & 0/3 & 0/3 & 0/3 \\
8 & 200 & 1 & 0\% & 1.23 & 3.36 & 4.32 & 0/3 & 0/3 & 1/3 \\
\midrule
\multicolumn{7}{@{}l}{\textbf{Total}} & \textbf{8/30} & \textbf{7/30} & \textbf{16/30} \\
\bottomrule
\end{tabular}

  \caption{\phynetpy\ and PhyloNet 3.8.2 running pseudo-likelihood network
  inference on byte-identical gene trees, each given a reticulation budget equal
  to the number of reticulations in the generating network. Three independently
  simulated replicates per condition; the right-hand block counts replicates in
  which the inferred network equalled the generating network ($\mu$-distance
  zero). ``Conc.'' is the mean percentage of simulated gene trees whose topology
  matches the generating network, so a low value is a harder problem. \phynetpy\
  is reported under its default search settings and under the \code{phylonet}
  preset, which reproduces PhyloNet's per-topology parameter optimisation.
  PhyloNet's runtime is self-reported and excludes JVM start-up.}
  \label{tab:phylonet}
\end{table}

\begin{figure}[t]
  \centering
  \includegraphics[width=\linewidth]{performance.pdf}
  \caption{(a) Runtime of pseudo-likelihood inference on identical input, on a
  log scale, one point per condition of Table~\ref{tab:phylonet}. Points are not
  joined because conditions sharing a taxon count differ in the number of gene
  trees. (b) Time to evaluate the exact biallelic-marker likelihood of a fixed
  network against the number of sites. (c) Dissimilarity-measure runtime against
  taxon count. Panels show runtime only; accuracy is in
  Table~\ref{tab:phylonet}. Every panel is drawn from the JSON written by the
  benchmark scripts.}
  \label{fig:performance}
\end{figure}

Gene trees were simulated under the multispecies network coalescent from a known
species tree or network, written once, and given to both implementations. Both
were asked for the same method -- pseudo-likelihood network inference, PhyloNet's
\code{InferNetwork_MPL} -- and both were given a reticulation budget equal to the
true count. The population parameter was set to $\theta = 0.5$, which leaves
between 46\% and none of the gene trees concordant with the generating topology
depending on the condition; at the simulator's smaller default scale every
simulated gene tree comes out identical to the species tree, and a comparison on
such data measures nothing. Both searches are stochastic, so each of the ten
conditions was run on three independently simulated replicates, giving 30 trials
per implementation. Table~\ref{tab:phylonet} and Figure~\ref{fig:performance}(a)
report the aggregate; per-replicate values are in the benchmark JSON.

On runtime, \phynetpy\ is faster on eight of the ten conditions, by up to about
$6\times$ at eight taxa with 200 gene trees, and the margin grows with problem
size. The two exceptions are the single-reticulation conditions at five and six
taxa, where PhyNetPy's mean runtime is the higher of the two. Because PhyloNet's
figure is self-reported and excludes JVM start-up, the comparison is conservative
in PhyloNet's favour. Both tools were run single-threaded.

The accuracy finding runs the other way, and it is the more consequential of the
two. PhyloNet recovered the generating topology exactly in 16 of 30 trials.
\phynetpy\ did so in 8 of 30 under its default settings and 7 of 30 under the
\code{phylonet} preset -- roughly half as often. The preset, which exists
precisely to match PhyloNet's per-topology parameter optimisation, does not close
the gap, and on one condition it costs several times the runtime for no
improvement. On this benchmark \phynetpy's pseudo-likelihood search is faster and
less accurate than PhyloNet's, and neither tool recovers the truth reliably on
the reticulate conditions.

We have not isolated the cause, and it would be premature to attribute it to the
search strategy. One observation argues against the simplest explanation, that
the search is merely under-run: on an eight-taxon condition it reports a log
pseudo-likelihood about two orders of magnitude worse than the score of the
generating network itself ($-9.2 \times 10^5$ against $-8.3 \times 10^3$), and
that value does not change when the iteration budget is raised from its default
to 50{,}000. More iterations would help a search that was stopping early. Two
candidate explanations remain open. Continuous parameters on the returned network
may be left un-optimised, so a near-correct topology is reported with a poor
score. And the pseudo-likelihood criterion does not enforce the branch-length
unit discipline that the full likelihood does, so a network carrying lengths on
the substitution scale is scored as though they were coalescent units. Both
are recorded as defects rather than explained away.\footnote{[TODO: resolve
before submission. Either fix the underlying cause and re-run this benchmark, or
keep this subsection as written and add the diagnosis. Do not soften the accuracy
result without new measurements.]}

Two limits on the comparison should be read with it. The simulated data come from
\phynetpy's own simulator under a model both implementations assume, which makes
the comparison reproducible from this repository and makes the runtime
measurement fair, but says nothing about robustness to model misspecification.
And while three replicates are enough to show that the accuracy difference is not
an artefact of a single lucky run, they are not enough to estimate its size
precisely; the counts in Table~\ref{tab:phylonet} should be read as indicating a
gap of roughly a factor of two rather than as a calibrated effect size.

\subsection{Component measurements}

\begin{table}[t]
  \centering
  \small
  % Generated by benchmarks/make_artifacts.py -- do not edit by hand.
\begin{tabular}{llrrrr}
\toprule
Model & Implementation & \textsc{PhyNetPy} & \texttt{expm} & Ratio & $\max|\Delta|$ \\
 & & ($\mu$s) & ($\mu$s) & & \\
\midrule
JC & closed form & 1.68 & 24.75 & 14.7$\times$ & 1.1e-16 \\
F81 & closed form & 1.32 & 24.96 & 18.9$\times$ & 1.4e-17 \\
K81 & closed form & 1.45 & 25.95 & 18.0$\times$ & 1.1e-16 \\
K80 & Tamura-Nei & 2.03 & 24.70 & 12.1$\times$ & 1.1e-16 \\
HKY & Tamura-Nei & 2.06 & 23.74 & 11.5$\times$ & 2.8e-17 \\
TN93 & Tamura-Nei & 2.04 & 24.78 & 12.1$\times$ & 1.1e-16 \\
SYM & cached eigendecomposition & 2.67 & 23.91 & 8.9$\times$ & 2.1e-16 \\
GTR & cached eigendecomposition & 2.60 & 25.45 & 9.8$\times$ & 4.4e-16 \\
\bottomrule
\end{tabular}

  \caption{Transition-matrix evaluation at $t = 0.1$ against a fresh
  \code{scipy.linalg.expm} call on the same rate matrix. Six models have a
  closed form; SYM and GTR cache an eigendecomposition. The final column is the
  largest absolute elementwise disagreement with \code{expm}, so this table is
  simultaneously a speed and a correctness measurement.}
  \label{tab:substitution}
\end{table}

\begin{table}[t]
  \centering
  \small
  % Generated by benchmarks/make_artifacts.py -- do not edit by hand.
\begin{tabular}{lrrr}
\toprule
Operation & \textsc{PhyNetPy} & NetworkX & Ratio \\
 & ($\mu$s) & ($\mu$s) & (NetworkX\,/\,\textsc{PhyNetPy}) \\
\midrule
Incident in-edges of a node & 0.16 & 1.81 & 11.45$\times$ \\
Out-degree of a node & 0.16 & 0.20 & 1.21$\times$ \\
Enumerate all leaves & 173.41 & 453.99 & 2.62$\times$ \\
Topological order & 2384.32 & 1295.26 & 0.54$\times$ \\
Leaf descendants, all nodes & 1452.37 & --- & --- \\
\midrule
Peak construction memory (KiB) & 2103 & 1111 & 0.53$\times$ \\
\bottomrule
\end{tabular}

  \caption{The compiled graph core against \code{networkx.DiGraph}
  \cite{Hagberg2008NetworkX} on the same 1{,}999-node topology. Ratios above
  one favour \phynetpy. \phynetpy\ is faster on the incidence and leaf queries
  that dominate network search and slower on topological ordering, and it uses
  roughly twice the peak memory during construction.}
  \label{tab:graphcore}
\end{table}

\begin{table}[t]
  \centering
  \small
  % Generated by benchmarks/make_artifacts.py -- do not edit by hand.
\begin{tabular}{rrrrrr}
\toprule
Taxa & Genes & Compiled (ms) & Python (ms) & Ratio & $|\Delta \log \mathrm{PL}|$ \\
\midrule
6 & 40 & 1.13 & 1.46 & 1.29$\times$ & 1.1e-13 \\
8 & 60 & 2.97 & 4.69 & 1.58$\times$ & 1.1e-12 \\
10 & 80 & 7.10 & 11.08 & 1.56$\times$ & 9.1e-13 \\
12 & 100 & 12.93 & 21.70 & 1.68$\times$ & 2.9e-11 \\
\bottomrule
\end{tabular}

  \caption{The compiled pseudo-likelihood dynamic program against the
  pure-Python reference that remains in the repository, selected by a module
  flag so that only the inner loop differs. The final column confirms the two
  paths return the same score.}
  \label{tab:kernels}
\end{table}

Three component measurements explain where \phynetpy's time goes.

Transition-matrix evaluation (Table~\ref{tab:substitution}) is between
$8.5\times$ and $19\times$ faster than calling \code{expm}, because six of the
eight models use a closed form and the remaining two cache an
eigendecomposition. This is not a claim to have improved on SciPy; it is the
expected benefit of exploiting structure that a general matrix exponential
cannot see.

The compiled graph core (Table~\ref{tab:graphcore}) is a genuinely mixed
result. Incidence lookup is about $11\times$ faster than NetworkX and
leaf enumeration about $3.8\times$, both essentially independent of network
size, which is the payoff from maintaining incidence indices. Topological
ordering is about $0.7\times$ -- that is, slower. And peak memory during
construction is roughly $1.9\times$ NetworkX's. The operations \phynetpy\
optimises are the ones a search loop performs most often, and the trade it makes
is memory for lookup speed; a reader whose workload is dominated by whole-graph
traversals rather than incidence queries should not expect a gain.

The compiled pseudo-likelihood kernel (Table~\ref{tab:kernels}) is worth
between $1.3\times$ and $1.7\times$ over the pure-Python dynamic program, with
the advantage growing with taxon count, and the two paths agree to within
$2 \times 10^{-12}$. That is a smaller factor than compiled-versus-interpreted
comparisons often report, because most of the work is already in array
operations rather than Python-level loops.

Biallelic-marker likelihood evaluation (Table~\ref{tab:snp},
Figure~\ref{fig:performance}(b)) takes about $1.1$\,ms for three taxa and
1{,}000 sites, scaling roughly linearly in sites and more steeply in taxa, as
the exact algorithm's dependence on the number of sampled lineages implies. The
first evaluation in a process costs noticeably more than subsequent ones
because the likelihood model is constructed on first use; both figures are
recorded, since a single interactive score and a score inside a search loop are
different quantities.

\begin{table}[t]
  \centering
  \small
  % Generated by benchmarks/make_artifacts.py -- do not edit by hand.
\begin{tabular}{rrrr}
\toprule
Taxa & Sites & Time per score (ms) & $\log L$ \\
\midrule
3 & 100 & 0.37 & -106.2 \\
3 & 1000 & 1.20 & -1030.6 \\
3 & 5000 & 4.94 & -5161.8 \\
4 & 100 & 0.69 & -114.2 \\
4 & 1000 & 1.73 & -1067.0 \\
4 & 5000 & 7.94 & -5774.7 \\
5 & 100 & 0.70 & -120.9 \\
5 & 1000 & 2.40 & -1196.6 \\
5 & 5000 & 11.89 & -6461.3 \\
\bottomrule
\end{tabular}

  \caption{Exact biallelic-marker likelihood of a fixed network, median of
  repeated evaluations with the likelihood model already constructed. The
  single-evaluation cost, which includes that construction, is recorded
  alongside these figures in \code{benchmarks/snp.json}.}
  \label{tab:snp}
\end{table}

\subsection{What is not measured}

Three gaps should be read alongside the tables above. No comparison against
PhyloNet's Bayesian samplers is reported: matching two MCMC implementations
requires agreeing on chain length, proposal weights and convergence criteria
before any timing is meaningful, and that comparison has not been done.
No empirical dataset is analysed here. And every measurement comes from one
machine, so the ratios should be read as indicative rather than as portable
constants.
